% 结论章节
% 说明：
% 1. 本文件设计为被主文档 \input{} 进来使用。

\section{结论}

本文提出了一种面向增量式 2DGS 场景重建的上下文感知前馈高斯回归系统，从表示与初始化两个层面改进增量式高斯建图的地图后端。在表示层面，系统以 2DGS surfel 替代传统 3DGS 体高斯，围绕 surfel 的法线语义构建了融合深度补全与单目法线先验的统一几何管线；在初始化层面，前馈回归网络通过跨注意力机制聚合多视角图像--渲染上下文，并以法线先验为锚点经 Gram-Schmidt 正交化构造 surfel 旋转，在单次前向传播中直接预测完整的 surfel 属性，替代传统基于规则的初始化。网络训练通过自蒸馏闭环完成，以高斯为核心单位构造训练样本，无需外部标注数据集。

在 R3LIVE 与 MCD 数据集上的实验表明，本文方法在全部测试序列上均取得了最优渲染指标，较最接近的基线 Gaussian-LIC2 在 R3LIVE 上平均提升 1.28\,dB PSNR。消融实验进一步表明，前馈回归器的核心价值在于改善新高斯插入后的初始质量与早期收敛速度：在固定 16 轮优化预算下，回归器使更多训练帧在有限优化步数内达到目标渲染质量，这对在线系统尤为关键——更好的初始高斯为后续帧提供更可靠的渲染监督，从而形成正向循环。

当前系统仍存在若干局限。首先，自蒸馏训练需要在目标场景上先以 naive 模式运行一遍完整建图以收集伪真值，增加了部署前的准备成本。其次，前馈网络在训练序列与测试序列之间的跨场景泛化能力尚未得到系统验证。未来工作将探索跨场景的泛化训练策略，并研究将前馈回归器与前端位姿估计更紧密耦合的可能性。
